
\sourcenote{Genre: core rule (tier 1) --- an ordinary mathematical proof
at Putnam-exam-solution detail. This rule concludes a \emph{satisfaction
judgement} ($\vDash$), one theory-extension edge above the base calculus:
every premise row is an inherited base-calculus sequent, and soundness
is proved \emph{given} base soundness (the extension-edge convention).
The theorems quoted above warrant the rule at the
\textbf{discrete-kernel profile}; the profile remark at the end of this
section states exactly what remains a hypothesis beyond it. Lemmas
\ref{lem:qsb-stopped}--\ref{lem:qsb-ville} and
Proposition~\ref{prop:qsb-pattern} are the battery family\textquotesingle s
shared probabilistic preamble.}

\textbf{Setting and claim.}\enspace Fix a measurable state space $S$ ---
the flattened hybrid state space, i.e.\ the disjoint sum of the per-mode
domains --- and a Markov one-step kernel $K$ on $S$, carrying both
discretized flow steps and guarded jumps. For $x \in S$ write $\Pr_x$
for the path measure on $S^{\N}$: the trajectory law of the
time-homogeneous chain driven by $K$, started deterministically at $x$
(the Ionescu--Tulcea construction, with the constant kernel family).
$(\mathcal{F}_n)_n$ is the canonical filtration, $\mathcal{F}_n$
generated by the coordinates $\le n$.

The rule\textquotesingle s data: a modelled domain $D$, a hull
$\bar{D} \supseteq D$ (the region on which the polynomial rows were
actually checked), an avoid set $U$ compiled to over-approximate the
falsity set of the $\mathsf{G}$-argument, a measurable certificate
$V : S \to \R$, and rationals $p, \mathrm{thr}$. The checked rows and
side conditions, read semantically at this profile:
\begin{itemize}
\item[(B1$\wedge$B2)] \emph{drift}: for every $y \in D \smallsetminus U$,
  $V$ is integrable against $K\,y$ and $\int V \,dK\,y \le V(y)$. (At
  this profile the interior and guard rows arrive \emph{together}: the
  one-step kernel carries flow and enabled jumps in a single step, so
  the two row families compile to one kernel-level drift row.)
\item[(B3)] \emph{floors}: $V \ge 1$ on $U \cap \bar{D}$.
\item[(B4)] \emph{nonnegs}: $V \ge 0$ on $\bar{D}$.
\item[(R3)] \emph{hull}: for every $y \in D \smallsetminus U$,
  $K\,y\,(\bar{D}) = 1$ --- one-step mass stays where the rows were
  checked.
\item[(R4)] \emph{fail-closed containment}: for every
  $y \in D \smallsetminus U$, $K\,y\,(D \cup U) = 1$.
\item[($\Theta$)] \emph{threshold}: $p + \mathrm{thr} \le 1$, checked in
  exact rational arithmetic.
\item[(F)] \emph{formula tie}: $\lnot \mathrm{sat}(y) \Rightarrow y \in U$,
  where $\mathrm{sat}$ is the state-satisfaction predicate of the
  $\mathsf{G}$-argument $\varphi$.
\end{itemize}
The claim (\texttt{sound\_QSafetyBattery}): for every $x \in D$ with
$V(x) \le \mathrm{thr}$,
\[
\Pr_x\bigl[\forall n,\ \mathrm{sat}(\omega_n)\bigr] \;\ge\; p .
\]

\begin{blemma}[the stopped certificate]\label{lem:qsb-stopped}
Let $f : S \to \R$ and let $A \subseteq S$ be a stop set. Define the
stopped certificate process $X_n(\omega) := f(\omega_{\,n \wedge \tau_A})$,
where $\tau_A$ is the first visit to $A$ --- equivalently, by the
primitive recursion that freezes the value once $A$ has been visited.
Then: (i) if $A$ is measurable and $f$ is measurable, $X$ is adapted to
$(\mathcal{F}_n)$; (ii) any pointwise property of $f$\textquotesingle s
values holds of $X$\textquotesingle s values; (iii) \emph{barrier
transport}: if $f \ge 1$ on $A$, then $X_n \ge 1$ on the event
$\{\tau_A \le n\}$; (iv) $X_n$ depends only on the coordinates $\le n$.
\end{blemma}
\begin{proof}
All four are elementary inductions on $n$ over the two branches of the
recursion (surviving: $X_{n+1} = f(\omega_{n+1})$; frozen:
$X_{n+1} = X_n$). No stopping-time API is needed anywhere --- the
recursion \emph{is} the stopped process, and (iii) is proved by chasing
the first index at which $A$ was visited.
\end{proof}

\begin{blemma}[supermartingale bridge]\label{lem:qsb-bridge}
Let $f$ be measurable with $0 \le f \le C$, let $A$ be measurable, and
suppose the one-step drift row holds \emph{off} $A$: for all
$y \notin A$, $\int f\,dK\,y \le f(y)$. Then for every start $x$, the
stopped certificate $X$ of Lemma~\ref{lem:qsb-stopped} is a
supermartingale with respect to $(\mathcal{F}_n, \Pr_x)$.
\end{blemma}
\begin{proof}
Integrability is the bound $0 \le X_n \le C$ (Lemma
\ref{lem:qsb-stopped}(ii)). For the supermartingale inequality,
condition on $\mathcal{F}_n$: by the Markov property of the trajectory
construction, the conditional expectation of a bounded
$\mathcal{F}_{n+1}$-measurable function given $\mathcal{F}_n$ is its
integral against the trajectory kernel continued from the $n$-prefix.
Split on survival at time $n$. On the surviving branch, paths patched to
agree with $\omega$ up to $n$ still survive
(Lemma~\ref{lem:qsb-stopped}(iv)), so the integrand is
$f(\text{next coordinate})$, whose one-step marginal is $K(\omega_n)$;
since survival means $\omega_n \notin A$, the drift row applies and the
conditional expectation is $\le f(\omega_n) = X_n(\omega)$. On the
frozen branch the integrand is the constant $X_n(\omega)$. Note that
\emph{every} hypothesis here is pointwise --- there is no
almost-everywhere invariance argument anywhere in the bridge; this is
what the fail-closed design buys, and the reachability sibling
(\rulename{QDeadlineBattery}) will have to pay precisely this cost when
it can no longer enlarge its stop set.
\end{proof}

\begin{blemma}[Ville\textquotesingle s inequality, barrier form]
\label{lem:qsb-ville}
Let $X$ be a nonnegative supermartingale with $X_n \ge 1$ on the event
$\{$some coordinate $\le n$ lies in $A\}$. Then
$\Pr_x[\exists n,\ \omega_n \in A] \le \mathbb{E}\,X_0$.
\end{blemma}
\begin{proof}
This is the maximal inequality for nonnegative supermartingales at
barrier $1$ \cite{ville1939} (textbook treatment: \cite{doob1953}): the
probability that a nonnegative supermartingale ever reaches level $a$ is
at most $\mathbb{E}\,X_0 / a$. The mechanized form deliberately asks no
measurability of the visit event itself --- the bad event is bounded by
inclusion in $\{\exists n,\ X_n \ge 1\}$, and only the process needs to
be adapted.
\end{proof}

\begin{bprop}[\textbf{the battery pattern}]\label{prop:qsb-pattern}
Let $A$ be measurable and let $f$ be measurable with
$0 \le f \le C$ globally, $f \ge 1$ on $A$, and one-step drift off $A$.
Then for every $x$,
\[
\Pr_x\bigl[\exists n,\ \omega_n \in A\bigr] \;\le\; f(x).
\]
\end{bprop}
\begin{proof}
Lemmas \ref{lem:qsb-stopped}--\ref{lem:qsb-ville} assembled; the initial
value is $\mathbb{E}\,X_0 = f(x)$ because coordinate $0$ is
deterministically $x$. This proposition --- \emph{nonnegative
supermartingale $+$ barrier $+$ initial value $\Rightarrow$ probability
bound} --- is the shape every rule in the battery family consumes; the
siblings cite it as ``the battery pattern
(\rulename{QSafetyBattery}, Proposition~\ref{prop:qsb-pattern})''.
\end{proof}

\begin{blemma}[the clamp: from checked rows to global hypotheses]
\label{lem:qsb-clamp}
The rows (B3), (B4) speak only on $\bar{D}$, while
Proposition~\ref{prop:qsb-pattern} wants its bounds globally. Set
$A := U \cup D^{\mathrm{c}}$ (the fail-closed stop set),
$C := \max(\mathrm{thr}, 1)$, and let $g$ be the clamped certificate:
on $D$, clamp $V$ into $[0, C]$; off $D$, clamp $\max(V, 1)$ into
$[0,C]$ (so the floor is \emph{forced} off the modelled domain). Then:
(i) $0 \le g \le C$ globally; (ii) $g \ge 1$ on $A$; (iii) on
$\bar{D} \cap (D \cup U)$ the clamp is invisible:
$g = \min(V, C)$ there; (iv) under (R3) $+$ (R4), $g$ satisfies the
one-step drift row off $A$; (v) $g(x) = V(x)$ at any
$x \in D$ with $V(x) \le \mathrm{thr}$.
\end{blemma}
\begin{proof}
(i), (ii), (v) are case checks from (B3), (B4) and the definition ---
(ii) uses (B3) on $U \cap \bar D$ and the forced floor off $D$. (iii):
on $D$ the lower clamp is inert by (B4); on
$(\bar{D} \cap U) \smallsetminus D$ the forced floor is inert by (B3).
For (iv), fix $y \in D \smallsetminus U$ (i.e.\ $y \notin A$). By (R3)
and (R4) the one-step measure $K\,y$ is carried by
$\bar{D} \cap (D \cup U)$ --- the complement has measure
$\le K y (\bar D^{\mathrm c}) + K y ((D \cup U)^{\mathrm c}) = 0$ ---
so by (iii) the integrand may be replaced by $\min(V, C)$. Then
$\int \min(V, C)\,dK\,y \le \min\bigl(\int V\,dK\,y,\; C\bigr)
\le \min(V(y), C) = g(y)$, using monotonicity of the integral, the
drift row (B1$\wedge$B2), and (iii) at $y$ itself.
\end{proof}

\begin{bthm}[fail-closed form]\label{thm:qsb-failclosed}
Under (B1)--(B4), (R3), (R4), ($\Theta$): for every $x \in D$ with
$V(x) \le \mathrm{thr}$,
\[
\Pr_x\bigl[\forall n,\ \omega_n \in D \smallsetminus U\bigr] \;\ge\; p .
\]
\end{bthm}
\begin{proof}
Apply Proposition~\ref{prop:qsb-pattern} to $(g, A)$ of
Lemma~\ref{lem:qsb-clamp}:
$\Pr_x[\text{reach } A] \le g(x) = V(x) \le \mathrm{thr}$. The
complement of ever visiting $A = U \cup D^{\mathrm{c}}$ is exactly
$\{\forall n,\ \omega_n \in D \smallsetminus U\}$, so its probability is
at least $1 - \mathrm{thr} \ge p$ by ($\Theta$). Note the concluded
event is \emph{stronger} than the safety formula\textquotesingle s truth
set: the chain never fails \emph{and never leaves the modelled domain}.
Leaving counts as failure --- that is the fail-closed reading, and it is
what (R4) is recomputed for.
\end{proof}

\begin{bprop}[soundness of \rulename{QSafetyBattery}]\label{prop:qsb-sound}
Under the hypotheses above including the formula tie (F): for every
$x \in D$ with $V(x) \le \mathrm{thr}$,
$\Pr_x[\forall n,\ \mathrm{sat}(\omega_n)] \ge p$ --- the concluded
satisfaction judgement
$\mathcal{M} \vDash_{\{V \le \mathrm{thr}\}}
\mathsf{P}_{\ge p}[\mathsf{G}\,\varphi]$, read on the sublevel region.
\end{bprop}
\begin{proof}
By the contrapositive of (F), $\omega_n \in D \smallsetminus U$ implies
$\mathrm{sat}(\omega_n)$; so the event of
Theorem~\ref{thm:qsb-failclosed} is contained in
$\{\forall n,\ \mathrm{sat}(\omega_n)\}$, and monotonicity of the
(outer) measure finishes. No measurability is ever asked of the
satisfaction predicate --- the smaller, fail-closed event is the
measurable one.
\end{proof}

\begin{bremark}[why the floors are not decoration]\label{rem:qsb-floors}
Drop (B3) and take $V \equiv 0$: drift, nonneg, and threshold rows all
hold, with $V(x) = 0 \le \mathrm{thr}$ for any $\mathrm{thr} \ge 0$, so
with $\mathrm{thr} = 0$, ($\Theta$) permits $p = 1$: the rule would
certify \emph{almost-sure} safety --- of any chain at all, including the
two-state chain $D = \{d\}$, $U = \{u\}$, $K\,d = \delta_u$, which fails
at its first step. With the floor present, this certificate dies
immediately: (B3) forces $V(u) \ge 1$, and then the drift row at $d$
demands $\int V\,dK\,d = V(u) \le V(d) = 0$ --- contradiction. Failure
must cost at least one unit of certificate value; the floor is the
price list.
\end{bremark}

\begin{bremark}[why the nonnegs are not decoration]\label{rem:qsb-nonnegs}
The maximal inequality of Lemma~\ref{lem:qsb-ville} is \emph{false} for
signed supermartingales, and the classical counterexample is worth
keeping in view. Take the simple symmetric random walk on $\mathbb{Z}$ started
at $0$, $A := \{1\}$, and the certificate $V(y) = y$. The drift row off
$A$ holds with equality (a martingale); the floor on $A$ holds
($V(1) = 1$); the initial value is $0$. Were
Proposition~\ref{prop:qsb-pattern} available without nonnegativity, it
would conclude $\Pr[\text{reach } 1] \le 0$ --- but the walk reaches $1$
almost surely. The signed certificate pays for its barrier visit with
mass borrowed from arbitrarily far below zero; (B4) is what forbids the
loan. (Boundedness plays only a technical integrability role; the
counterexample targets the sign.)
\end{bremark}

\begin{bremark}[the hull and fail-closed conditions (R3), (R4)]
\label{rem:qsb-r34}
Both are recomputed by the checker, never taken from the producer, and
each blocks a distinct hole. Without (R3), the one-step integral in the
drift row could consult states outside $\bar{D}$ --- states where the
floor and nonneg rows were \emph{never checked} --- and
Lemma~\ref{lem:qsb-clamp}(iv)\textquotesingle s clamp-invisibility
computation collapses. Without (R4), one-step mass could leave
$D \cup U$ entirely: the concluded event counts leaving the modelled
domain as failure, and (R4) is exactly what makes the clamped
certificate\textquotesingle s drift hold at the \emph{enlarged} stop set
$U \cup D^{\mathrm{c}}$ --- the design that let every hypothesis in
Lemma~\ref{lem:qsb-bridge} stay pointwise. A companion condition
(floor-scope resolution) is discharged at row-compilation time: it is
what makes the checked floor rows imply (B3) as stated here, and it has
no semantic residue at this layer.
\end{bremark}

\begin{bremark}[threshold and integrability]\label{rem:qsb-theta}
($\Theta$) is bookkeeping, and it is sharp: the certificate buys
$\Pr[\mathrm{fail}] \le \mathrm{thr}$, so asserting success
$\ge p$ requires $p \le 1 - \mathrm{thr}$, and
$\mathrm{thr} = 1 - p$ is the tightest admissible choice --- producer
weakening ($\mathrm{thr} < 1 - p$) is sound for free. The integrability
half of (B1$\wedge$B2) is the discrete-profile reading of
``the certificate lies in the generator\textquotesingle s
domain''; for the intended instances --- polynomial $V$, one-step
images inside the compact hull by (R3) --- it is automatic.
\end{bremark}

\begin{bremark}[\textbf{profile} --- what is proven, what is filed]
\label{rem:qsb-profile}
Proven, and quoted above this proof: the discrete-kernel profile ---
one-step Markov kernel semantics over the trajectory construction, the
theorems \texttt{sound\_QSafetyBattery} and
\texttt{sound\_QSafetyBattery\_failClosed}, sorry-free. Also proven,
end-to-end: the \emph{PDMP sub-profile} of the continuous-time semantics
--- piecewise-deterministic flow with guarded jumps, no diffusion
(\texttt{sound\_GSHS\_safety\_battery}). \emph{Not} proven: the general
continuous-time profile (jump--diffusions), where this same argument
must ride continuous-time generator domains, fine properties of sample
paths, and optional stopping for right-continuous supermartingales. At
that profile the rule\textquotesingle s warrant is a stated hypothesis,
and the present section is testimony to the shape of the argument, not a
theorem. The conditional mark on the dashboard row records exactly this
gap --- nothing in this book should be read as narrowing it.
\end{bremark}

\begin{bremark}[the staged premise shapes and the system bridge]
\label{rem:qsb-staged}
The rule\textquotesingle s node presents its premises per mode and per
jump: mode-indexed certificate family $(V_m)$, interior rows against
each mode\textquotesingle s flow mechanism, guard rows per jump, floors
and nonnegs per mode, and the threshold row in exact rational
arithmetic, cast once. The bound theorem for that shape
(\texttt{sound\_QSafetyBattery\_staged}) closes the gap to the flattened
hypotheses used above through an explicit bridge hypothesis: at every
in-scope state, the one-step kernel decomposes as a sub-convex mixture
of the mode\textquotesingle s flow step and the enabled guard-reset
jumps (weights supported on enabled jumps). Under that decomposition the
mode\textquotesingle s interior row and the enabled guard rows
\emph{are} the flattened drift row (\texttt{mixture\_drift}): the
one-step integral splits over the mixture, each jump term is a guard row
evaluation, and the weights sum to one. What stands left of the
concluded $\vDash$ is any system presentation whose compilation at this
profile satisfies the bridge.
\end{bremark}

\begin{bremark}[relation to the literature]\label{rem:qsb-lit}
Supermartingale safety certificates for stochastic systems go back to
Kushner \cite{kushner1967}; the modern certificate-form idiom ---
synthesize a barrier function, check polynomial rows, conclude a
probability bound --- is Prajna--Jadbabaie--Pappas
\cite{prajnajadbabaiepappas2007}. The probability engine is
Ville\textquotesingle s inequality \cite{ville1939}, with the systematic
supermartingale treatment in Doob \cite{doob1953}. The reading of the
interior rows as one-step expected decrease is the Foster--Lyapunov
drift idiom \cite{meyntweedie2009}, whose geometric form powers the
deadline sibling. What this rule adds to that literature is not
mathematics but \emph{trust structure}: every analytic ingredient lives
in a verified lemma, every polynomial ingredient is an inherited
base-calculus row replayed by the kernel, and the checker recomputes the
side conditions rather than trusting the producer.
\end{bremark}

\begin{anchors}
Lemma \ref{lem:qsb-stopped}: stopped certificate, adaptedness, barrier transport & \texttt{stopped}, \texttt{survive}, \texttt{stronglyAdapted\_stopped}, \texttt{one\_le\_stopped}, \texttt{stopped\_congr}\\
Lemma \ref{lem:qsb-bridge}: supermartingale bridge & \texttt{supermartingale\_stopped} (Markov property via \texttt{Kernel.condExp\_traj}; prefix patching via \texttt{Kernel.integral\_traj}; one-step marginal via \texttt{Kernel.map\_traj\_succ\_self})\\
Lemma \ref{lem:qsb-ville}: maximal inequality at barrier $1$ & \texttt{Qtsl.Ville.reach\_avoid}, consumed at $a = 1$\\
Proposition \ref{prop:qsb-pattern}: the battery pattern & \texttt{measure\_reach\_le}\\
Lemma \ref{lem:qsb-clamp}: the clamp & the construction $g$, \texttt{hclamp}, \texttt{hgdrift} inside \texttt{sound\_QSafetyBattery\_failClosed}\\
Theorem \ref{thm:qsb-failclosed} & \texttt{sound\_QSafetyBattery\_failClosed}\\
Proposition \ref{prop:qsb-sound}: formula tie, monotonicity & \texttt{sound\_QSafetyBattery}\\
Remark \ref{rem:qsb-staged}: staged shapes, mixture bridge & \texttt{mixture\_integrable}, \texttt{mixture\_drift}, \texttt{sound\_QSafetyBattery\_staged}\\
Remark \ref{rem:qsb-profile}: PDMP sub-profile & \texttt{sound\_GSHS\_safety\_battery}\\
\end{anchors}
